\documentclass[11pt]{article}

\usepackage[utf8]{inputenc}
\usepackage[T1]{fontenc}
\usepackage{lmodern}
\usepackage{geometry}
\usepackage{amsmath,amssymb,amsfonts,amsthm,mathtools}
\usepackage{bm}
\usepackage{enumitem}
\usepackage{microtype}
\usepackage{hyperref}
\usepackage{titlesec}
\usepackage{booktabs}
\usepackage{array}
\usepackage{setspace}
\usepackage{mathrsfs}
\usepackage{physics}

\geometry{margin=1in}
\hypersetup{colorlinks=true, linkcolor=black, urlcolor=blue, citecolor=black}
\setlist[enumerate]{topsep=0.4em,itemsep=0.35em}
\setlist[itemize]{topsep=0.3em,itemsep=0.25em}
\titleformat{\section}{\large\bfseries}{\thesection.}{0.5em}{}
\titleformat{\subsection}{\normalsize\bfseries}{\thesubsection.}{0.5em}{}
\setstretch{1.06}

\newcommand{\Aether}{\AE{}ther}
\newcommand{\Aflow}{\AE{}ther-flow}
\newcommand{\TheoryName}{\Aflow{} Interpretation of Relativity}
\newcommand{\TheoryProgram}{\Aether{} / \Aflow{} framework}
\newcommand{\Sub}{\mathcal A}
\newcommand{\ObserverMap}{\Xi}
\newcommand{\BridgeMetric}{\mathfrak g}
\newcommand{\ObserverMetric}{g^{\mathrm{br}}}
\newcommand{\CandidateMetric}{g^{\mathrm{cand}}}
\newcommand{\CompletedMetric}{g^{\sharp}}
\newcommand{\BaseShift}{\Sigma^{\mathrm{IER}}}
\newcommand{\ResponseShift}{\Sigma^{\mathrm{resp}}}
\newcommand{\CompShift}{\Sigma^{\mathrm{cmp}}}
\newcommand{\BaseLift}{\widehat\Sigma^{\mathrm{IER}}}
\newcommand{\ResponseLift}{\widehat\Sigma^{\mathrm{resp}}}
\newcommand{\CompLift}{\widehat\Sigma^{\mathrm{cmp}}}
\newcommand{\ReLongFour}{\widetilde{\mathcal R}_{\mu\nu}^{(\chi,\ge 4)}}
\newcommand{\ResidualCore}{\mathcal V_{\mu\nu}^{\mathrm{res}}}
\newcommand{\ResidualLift}{\widehat{\mathcal V}^{\mathrm{res}}}
\newcommand{\SubEinLin}{\widehat{\mathcal E}}
\newcommand{\SubGaugeVec}{\widehat{\mathcal B}}
\newcommand{\SubGaugeBook}{\widehat{\mathcal G}}
\newcommand{\SubReducedEin}{\widehat{\mathcal L}}
\newcommand{\FreeProj}{\Pi_{\mathrm{free}}}
\newcommand{\GaugeAux}{Y}
\newcommand{\ReducedEin}{\mathcal L_{\ObserverMetric}}
\newcommand{\GaugeBook}{\mathcal G_{\ObserverMetric}}
\newcommand{\EinLin}{\mathcal E_{\ObserverMetric}}
\newcommand{\EinQuad}{\mathcal Q_{\ge 2}^{\mathrm{Ein}}}
\newcommand{\CompClass}{\mathscr C_{\mathrm{cmp}}}
\newcommand{\CompFree}{\mathscr C_{\mathrm{cmp}}^{\mathrm{free}}}
\newcommand{\CompForced}{\mathscr C_{\mathrm{cmp}}^{\mathrm{for}}}
\newcommand{\CompKernel}{\mathscr K_{\mathrm{cmp}}}
\newcommand{\MicFunctional}{\mathscr F_{\mathrm{mic}}}
\newcommand{\PrimFunctional}{\mathscr F_{\mathrm{prim}}}
\newcommand{\CarrierFunctional}{\mathscr F_{\mathrm{car}}}
\newcommand{\DeepFunctional}{\mathscr F_{\mathrm{deep}}}
\newcommand{\ProtoFunctional}{\mathscr F_{\mathrm{proto}}}
\newcommand{\UltraFunctional}{\mathscr F_{\mathrm{micro}}}
\newcommand{\RootFunctional}{\mathscr F_{\mathrm{hess}}}
\newcommand{\FreeStiff}{\kappa_{\mathrm{free}}}
\newcommand{\PrimStrain}{K}
\newcommand{\PrimNeutral}{J}
\newcommand{\PrimFree}{F}
\newcommand{\CarrierClass}{\mathscr H_{\mathrm{car}}}
\newcommand{\CarrierProj}{\Pi_{\mathrm{str}}}
\newcommand{\CarrierPerp}{\Pi_{\mathrm{str}}^{\perp}}
\newcommand{\CarrierGap}{\kappa_{\mathrm{str}}}
\newcommand{\CarrierSeed}{\mathfrak S}
\newcommand{\RelabelSeed}{\mathfrak R}
\newcommand{\FreeSeed}{\mathfrak F}
\newcommand{\PreCarGap}{\kappa_{\mathrm{car},0}}
\newcommand{\CurrentGap}{\kappa_{\mathrm{cur}}}
\newcommand{\HomGap}{\kappa_{\mathrm{hom}}}
\newcommand{\epsIR}{\varepsilon}
\newcommand{\DeepTensorClass}{\mathscr X_{\mathrm{car}}}
\newcommand{\DeepRelabelClass}{\mathscr X_{\mathrm{cur}}}
\newcommand{\DeepFreeClass}{\mathscr X_{\mathrm{hom}}}
\newcommand{\DeepTensor}{\mathfrak X}
\newcommand{\DeepRelabel}{\mathfrak J}
\newcommand{\DeepFree}{\mathfrak W}
\newcommand{\CarRead}{\mathcal R_{\mathrm{car}}}
\newcommand{\CurRead}{\mathcal R_{\mathrm{cur}}}
\newcommand{\HomRead}{\mathcal R_{\mathrm{hom}}}
\newcommand{\ProtoTensorClass}{\mathscr Y_{\mathrm{car}}}
\newcommand{\ProtoRelabelClass}{\mathscr Y_{\mathrm{cur}}}
\newcommand{\ProtoFreeClass}{\mathscr Y_{\mathrm{hom}}}
\newcommand{\ProtoTensor}{\mathfrak Y}
\newcommand{\ProtoRelabel}{\mathfrak K}
\newcommand{\ProtoFree}{\mathfrak Z}
\newcommand{\ProtoCarProj}{\Pi_{\mathrm{deep,car}}}
\newcommand{\ProtoCurProj}{\Pi_{\mathrm{deep,cur}}}
\newcommand{\ProtoHomProj}{\Pi_{\mathrm{deep,hom}}}
\newcommand{\ProtoCarPerp}{\Pi_{\mathrm{deep,car}}^{\perp}}
\newcommand{\ProtoCurPerp}{\Pi_{\mathrm{deep,cur}}^{\perp}}
\newcommand{\ProtoHomPerp}{\Pi_{\mathrm{deep,hom}}^{\perp}}
\newcommand{\ProtoCarRead}{\widetilde{\mathcal R}_{\mathrm{car}}}
\newcommand{\ProtoCurRead}{\widetilde{\mathcal R}_{\mathrm{cur}}}
\newcommand{\ProtoHomRead}{\widetilde{\mathcal R}_{\mathrm{hom}}}
\newcommand{\ProtoCarGap}{\gamma_{\mathrm{deep,car}}}
\newcommand{\ProtoCurGap}{\gamma_{\mathrm{deep,cur}}}
\newcommand{\ProtoHomGap}{\gamma_{\mathrm{deep,hom}}}
\newcommand{\UltraTensorClass}{\mathscr Z_{\mathrm{car}}}
\newcommand{\UltraRelabelClass}{\mathscr Z_{\mathrm{cur}}}
\newcommand{\UltraFreeClass}{\mathscr Z_{\mathrm{hom}}}
\newcommand{\UltraTensor}{\mathfrak M}
\newcommand{\UltraRelabel}{\mathfrak N}
\newcommand{\UltraFree}{\mathfrak P}
\newcommand{\UltraCarProj}{\Pi_{\mathrm{mic,car}}}
\newcommand{\UltraCurProj}{\Pi_{\mathrm{mic,cur}}}
\newcommand{\UltraHomProj}{\Pi_{\mathrm{mic,hom}}}
\newcommand{\UltraCarPerp}{\Pi_{\mathrm{mic,car}}^{\perp}}
\newcommand{\UltraCurPerp}{\Pi_{\mathrm{mic,cur}}^{\perp}}
\newcommand{\UltraHomPerp}{\Pi_{\mathrm{mic,hom}}^{\perp}}
\newcommand{\UltraCarGap}{\delta_{\mathrm{mic,car}}}
\newcommand{\UltraCurGap}{\delta_{\mathrm{mic,cur}}}
\newcommand{\UltraHomGap}{\delta_{\mathrm{mic,hom}}}
\newcommand{\TensorDefect}{\mathcal D_{\mathrm{car}}}
\newcommand{\CurDefect}{\mathcal D_{\mathrm{cur}}}
\newcommand{\HomDefect}{\mathcal D_{\mathrm{hom}}}
\newcommand{\UltraCarRead}{\mathcal S_{\mathrm{car}}}
\newcommand{\UltraCurRead}{\mathcal S_{\mathrm{cur}}}
\newcommand{\UltraHomRead}{\mathcal S_{\mathrm{hom}}}
\newcommand{\RootTensorClass}{\mathscr U_{\mathrm{car}}}
\newcommand{\RootRelabelClass}{\mathscr U_{\mathrm{cur}}}
\newcommand{\RootFreeClass}{\mathscr U_{\mathrm{hom}}}
\newcommand{\RootTensor}{\mathfrak U}
\newcommand{\RootRelabel}{\mathfrak V}
\newcommand{\RootFree}{\mathfrak Q}
\newcommand{\RootCarProj}{\Pi_{\mathrm{sel,car}}}
\newcommand{\RootCurProj}{\Pi_{\mathrm{sel,cur}}}
\newcommand{\RootHomProj}{\Pi_{\mathrm{sel,hom}}}
\newcommand{\RootCarPerp}{\Pi_{\mathrm{sel,car}}^{\perp}}
\newcommand{\RootCurPerp}{\Pi_{\mathrm{sel,cur}}^{\perp}}
\newcommand{\RootHomPerp}{\Pi_{\mathrm{sel,hom}}^{\perp}}
\newcommand{\TensorSymProj}{\Pi_{\mathrm{sym,car}}}
\newcommand{\CurSymProj}{\Pi_{\mathrm{sym,cur}}}
\newcommand{\HomSymProj}{\Pi_{\mathrm{sym,hom}}}
\newcommand{\TensorCompat}{\mathcal C_{\mathrm{car}}}
\newcommand{\CurCompat}{\mathcal C_{\mathrm{cur}}}
\newcommand{\HomCompat}{\mathcal C_{\mathrm{hom}}}
\newcommand{\TensorHess}{\mathcal H_{\mathrm{car}}}
\newcommand{\CurHess}{\mathcal H_{\mathrm{cur}}}
\newcommand{\HomHess}{\mathcal H_{\mathrm{hom}}}

\newtheorem{definition}{Definition}
\newtheorem{proposition}{Proposition}
\newtheorem{remark}{Remark}
\newtheorem{corollary}{Corollary}
\newtheorem{theorem}{Theorem}

\title{\textbf{The \TheoryName{}}\\[0.4em]
\large Substrate-Response / Self-Response Charge-Polarization Candidate\\
\large Quartic-Completion Selection-Principle Primitive-Reservoir Carrier-Origin\\
\large Precursor-Selection Exact-Preservation Normal-Form Microphysical-Hessian Origin Note}
\author{}
\date{}

\begin{document}

\maketitle

\begin{abstract}
The exact-preservation normal-form microphysical-origin note already shows that the ambient projected-sector construction is not primitive at present scope. It is the defect-null exact-preserving sector of microphysical spaces
\[
\UltraTensorClass,\qquad
\UltraRelabelClass,\qquad
\UltraFreeClass
\]
equipped with defect operators
\[
\TensorDefect,\qquad
\CurDefect,\qquad
\HomDefect,
\]
whose closed kernels recover the ambient spaces, whose orthogonal kernel projectors recover the admissible ambient sector, and whose positive-gap defect estimates suppress the exact-breaking complements. What remained open is one layer deeper again: why those microphysical spaces, those defect operators, their closed-kernel/orthogonal-projector structure, and their coercive positive-gap estimates should exist at all rather than being the present deepest disciplined postulate.

The present note addresses exactly that burden at sharper symmetry-compatible equilibrium scope. It introduces still sharper substrate spaces
\[
\RootTensorClass,\qquad
\RootRelabelClass,\qquad
\RootFreeClass
\]
with exact-preserving symmetry projectors
\[
\TensorSymProj,\qquad
\CurSymProj,\qquad
\HomSymProj
\]
and closed compatibility operators
\[
\TensorCompat,\qquad
\CurCompat,\qquad
\HomCompat.
\]
The admissible microphysical spaces are then forced to be the closed symmetry-fixed compatibility sectors
\[
\UltraTensorClass=\operatorname{ran}\TensorSymProj\cap\ker\TensorCompat,\qquad
\UltraRelabelClass=\operatorname{ran}\CurSymProj\cap\ker\CurCompat,\qquad
\UltraFreeClass=\operatorname{ran}\HomSymProj\cap\ker\HomCompat.
\]
On those induced microphysical sectors, the quadratic exact-preservation response is represented by nonnegative self-adjoint reduced Hessians
\[
\TensorHess,\qquad
\CurHess,\qquad
\HomHess.
\]
Their square roots are exactly the defect operators of the previous note:
\[
\TensorDefect=\TensorHess^{1/2},\qquad
\CurDefect=\CurHess^{1/2},\qquad
\HomDefect=\HomHess^{1/2}.
\]
The exact-preserving ambient sectors are therefore the zero-response sectors
\[
\ProtoTensorClass=\ker\TensorHess=\ker\TensorDefect,\qquad
\ProtoRelabelClass=\ker\CurHess=\ker\CurDefect,\qquad
\ProtoFreeClass=\ker\HomHess=\ker\HomDefect.
\]
Because the reduced Hessians are nonnegative self-adjoint operators on Hilbert spaces, those kernels are closed and the orthogonal kernel projectors are forced. If the positive parts of the Hessians carry spectral gaps
\[
\langle \TensorHess\UltraTensor,\UltraTensor\rangle
\ge
\UltraCarGap\norm{\UltraCarPerp\UltraTensor}_{\mathrm{micro,car}}^2,
\qquad
\langle \CurHess\UltraRelabel,\UltraRelabel\rangle
\ge
\UltraCurGap\norm{\UltraCurPerp\UltraRelabel}_{\mathrm{micro,cur}}^2,
\]
\[
\langle \HomHess\UltraFree,\UltraFree\rangle
\ge
\UltraHomGap\norm{\UltraHomPerp\UltraFree}_{\mathrm{micro,hom}}^2,
\qquad
\UltraCarGap,\UltraCurGap,\UltraHomGap>0,
\]
then the previous defect estimates are recovered verbatim.

The reduced symmetry-compatible equilibrium functional is therefore
\[
\RootFunctional[\UltraTensor,\UltraRelabel,\UltraFree]
\equiv
\ProtoFunctional[
\UltraCarProj\UltraTensor,\UltraCurProj\UltraRelabel,\UltraHomProj\UltraFree]
+\frac12\langle \TensorHess\UltraTensor,\UltraTensor\rangle
+\frac12\langle \CurHess\UltraRelabel,\UltraRelabel\rangle
+\frac12\langle \HomHess\UltraFree,\UltraFree\rangle,
\]
which is exactly the same microphysical exact-preservation functional already recorded one layer above:
\[
\RootFunctional[\UltraTensor,\UltraRelabel,\UltraFree]
=
\UltraFunctional[\UltraTensor,\UltraRelabel,\UltraFree].
\]
Consequently the same ambient projected-sector construction, the same exact-preservation normal form, the same carrier-origin functional, the same primitive reservoir, the same microscopic-equilibrium reservoir, the same \(\Pi_{\mathrm{free}}H=0\) outcome, the same completion solve \(\Sigma^{\mathrm{cmp}}\), the same observer-equation closure, and the same one operative metric are recovered unchanged.

The scope remains disciplined. This note does not classify every possible still sharper symmetry or compatibility structure, and it does not claim a final ultraviolet completion. Its exact positive result is narrower and precise: the defect-null microphysical construction of the previous note is itself forced by a sharper symmetry-compatible equilibrium-response structure, while preserving the same downstream one-metric completion package.
\end{abstract}

\section{Introduction}

The active charge-polarization line has now crossed the candidate, gate, controlled-limit, residual-sector, redesign, substrate-anchoring, substrate-action, kernel-origin, microscopic-equilibrium deeper-origin, equilibrium-reservoir primitive-origin, primitive-reservoir carrier-origin, carrier-origin precursor-selection, carrier-origin precursor-selection deeper-origin, exact-preservation normal-form origin, and exact-preservation normal-form microphysical-origin steps on the same one-metric GR route. The immediate burden therefore narrows again. It is not another packaging pass. It is not stronger promotion language. It is not, by default, another anchored positive-pair continuation. It is to go one layer below the defect-null microphysical construction now on record and explain why that construction itself is forced by still sharper substrate structure.

At current scope that burden has four linked parts. First, one must explain why the microphysical spaces
\[
\UltraTensorClass,\qquad
\UltraRelabelClass,\qquad
\UltraFreeClass
\]
are admissible at all rather than simply written down. Second, one must explain why the defect operators
\[
\TensorDefect,\qquad
\CurDefect,\qquad
\HomDefect
\]
are the correct exact-preservation defects rather than arbitrary quadratic penalties. Third, one must explain why their kernels are closed Hilbert subspaces with orthogonal kernel projectors and why the hidden defect-bearing complements carry coercive positive gaps. Fourth, one must re-show that once this sharper origin is imposed, the same downstream package survives without reopening the ambient projected-sector construction, the raw deeper readout spaces, or the observer-equation gains already recorded.

The present note addresses exactly those tasks. Its key move is to pass from the microphysical layer of the previous note to a sharper symmetry-compatible equilibrium layer. At that sharper level, the admissible microphysical sectors are not primitive spaces but the closed symmetry-fixed compatibility sectors of still sharper substrate spaces. The microphysical defect operators are then not primitive maps either: they are the square roots of the reduced equilibrium Hessians on those symmetry-compatible sectors. Their kernels are therefore the zero-response directions of the reduced response, and the orthogonal kernel projectors and coercive positive-gap estimates follow from standard Hilbert-space spectral structure. Once that reduced quadratic functional is written in Hessian form, it is exactly the same microphysical exact-preservation functional already used one layer above, so the same downstream elimination chain is recovered unchanged.

\paragraph{Framework and claim boundary.}
\TheoryProgram{} denotes the broader \Aether{} / \Aflow{} framework built on the claims that the \Aether{} is the underlying four-dimensional substrate of reality, the \Aflow{} is its intrinsic ordered motion, observed three-dimensional space is the local experiential slice of that deeper substrate, S-time is the experienced order of change arising from matter, light, and the \Aflow{}, and observed expansion is the three-dimensional appearance of deeper four-dimensional ordered motion. Within that framework, \TheoryName{} denotes the exact-closure theory in which the effective gravitational dynamics are adopted to be exactly Einsteinian with universal matter coupling. In the active sequence, \emph{adoption} means use of that established relativistic dynamics without claiming substrate derivation, while \emph{derivation} is reserved for a first-principles recovery from explicit substrate variables.

The present note stays strictly inside that boundary. It does not claim that every conceivable still sharper symmetry or compatibility principle has now been classified. It does not widen the theorem boundary beyond the current one-metric exact-preservation claim. It does make one narrower positive move: it shows that the previous defect-null microphysical construction is itself induced by a sharper symmetry-compatible equilibrium-response structure, while preserving the same carrier-origin functional, the same primitive reservoir, the same microscopic-equilibrium reservoir, the same \(\Pi_{\mathrm{free}}H=0\) outcome, the same \(\Sigma^{\mathrm{cmp}}\) solve, and the same one operative metric.

\section{Fixed Inputs from the Recorded Completion Line}

\begin{definition}[Fixed charge-polarization branch and source]
\label{def:fixed}
The present note takes as fixed the same charge-polarization branch
\begin{equation}
\CandidateMetric
=
\ObserverMetric+\BaseShift+\ResponseShift,
\qquad
\CompletedMetric
=
\ObserverMetric+\BaseShift+\ResponseShift+\CompShift,
\label{eq:metrics}
\end{equation}
with lifted completion variables \(\BaseLift,\ResponseLift,\CompLift\) satisfying
\begin{equation}
\ObserverMap^*\BaseLift=\BaseShift,
\qquad
\ObserverMap^*\ResponseLift=\ResponseShift,
\qquad
\ObserverMap^*\CompLift=\CompShift.
\label{eq:lifts}
\end{equation}
The unresolved observer-side quartic tensor and its fixed substrate lift are
\begin{equation}
\ResidualCore
\equiv
\ReLongFour
+\EinQuad[\BaseShift]
+\EinLin(\ResponseShift),
\label{eq:vres}
\end{equation}
\begin{equation}
\ResidualLift_{AB}
\equiv
\widehat{\mathcal R}_{AB}^{(\chi,\ge 4)}
+\widehat{\mathcal Q}_{AB}^{\mathrm{Ein}}[\BaseLift]
+\widehat{\mathcal E}_{AB}(\ResponseLift),
\qquad
\ObserverMap^*\ResidualLift=\ResidualCore.
\label{eq:liftedsource}
\end{equation}
The fixed orders are
\begin{equation}
\BaseShift=\mathcal O_{\ge 2},
\qquad
\ResponseShift=\mathcal O_{\ge 4},
\qquad
\CompShift=\mathcal O_{\ge 4},
\qquad
\ResidualLift=\mathcal O_{\ge 4},
\label{eq:orders}
\end{equation}
and on every controlled infrared family,
\begin{equation}
\BaseShift=\mathcal O(\epsIR^2),
\qquad
\ResponseShift=\mathcal O(\epsIR^4),
\qquad
\CompShift=\mathcal O(\epsIR^4),
\qquad
\ResidualLift=\mathcal O(\epsIR^4).
\label{eq:irorders}
\end{equation}
\end{definition}

\begin{definition}[Fixed bridge-response operators]
\label{def:operators}
Let \(\widehat\nabla\) denote the Levi-Civita connection of the unshifted bridge metric \(\BridgeMetric\). For any observer-compatible symmetric substrate tensor \(H_{AB}\), define
\begin{equation}
\SubGaugeVec(H)_A
\equiv
\widehat\nabla^B\!\left(
H_{AB}
-\frac12 \BridgeMetric_{AB}\BridgeMetric^{CD}H_{CD}
\right),
\label{eq:subB}
\end{equation}
\begin{equation}
\SubGaugeBook(H)_{AB}
\equiv
\widehat\nabla_{(A}\SubGaugeVec(H)_{B)}
-\frac12 \BridgeMetric_{AB}\widehat\nabla^C\SubGaugeVec(H)_C,
\label{eq:subG}
\end{equation}
\begin{equation}
\SubReducedEin(H)_{AB}
\equiv
\SubEinLin(H)_{AB}-\SubGaugeBook(H)_{AB},
\label{eq:subL}
\end{equation}
chosen so that
\begin{equation}
\ObserverMap^*\bigl(\SubEinLin(H)\bigr)=\EinLin(\ObserverMap^*H),
\qquad
\ObserverMap^*\bigl(\SubGaugeBook(H)\bigr)=\GaugeBook(\ObserverMap^*H),
\label{eq:intertwine1}
\end{equation}
and hence
\begin{equation}
\ObserverMap^*\bigl(\SubReducedEin(H)\bigr)
=
\ReducedEin(\ObserverMap^*H).
\label{eq:intertwine2}
\end{equation}
\end{definition}

\begin{definition}[Admissible completion class and free sector]
\label{def:class}
Let \(\CompClass\) denote the same observer-compatible reduced-harmonic Coulomb-plus-regular completion class used throughout the recorded charge-polarization line. The free homogeneous sector is
\begin{equation}
\CompFree
\equiv
\left\{
H\in\CompClass:\SubReducedEin(H)=0
\right\},
\label{eq:freeclass}
\end{equation}
and \(\CompForced\) denotes its orthogonal complement inside \(\CompClass\) with respect to the bridge-metric \(L^2\) pairing
\begin{equation}
\langle H_1,H_2\rangle_{\mathrm{cmp}}
\equiv
\int_{\Sub} d^4X\,\sqrt{G}\,
\BridgeMetric^{AC}\BridgeMetric^{BD}(H_1)_{AB}(H_2)_{CD}.
\label{eq:inner}
\end{equation}
Let \(\FreeProj:\CompClass\to\CompFree\) be the orthogonal projector. The fixed source satisfies
\begin{equation}
\ResidualLift\in\CompForced,
\qquad
\FreeProj\ResidualLift=0.
\label{eq:sourcefree}
\end{equation}
\end{definition}

\begin{definition}[Carrier-origin block fixed by the previous notes]
\label{def:carrier}
The primitive-reservoir carrier-origin note fixed a still-deeper carrier block with tensor seed
\[
\CarrierSeed\in\CarrierClass,
\qquad
\CarrierClass=\CompClass\oplus \CompClass^\perp,
\]
orthogonal projector
\[
\CarrierProj:\CarrierClass\to\CompClass,
\qquad
\CarrierPerp\equiv I-\CarrierProj,
\]
relabeling-current seed \(\RelabelSeed_A\in T^*\Sub\), and source-free homogeneous seed \(\FreeSeed_{AB}\in\CompFree\). Its functional is
\begin{equation}
\begin{aligned}
\CarrierFunctional[\CarrierSeed,\RelabelSeed,\FreeSeed]
\equiv
\int_{\Sub} d^4X\,\sqrt{G}\,
\Bigl[
&\frac12 (\CarrierProj\CarrierSeed)^{AB}
\SubEinLin(\CarrierProj\CarrierSeed)_{AB}
+\frac{\CarrierGap}{2}(\CarrierPerp\CarrierSeed)^{AB}(\CarrierPerp\CarrierSeed)_{AB}
\\
&+\frac12\bigl(\RelabelSeed_A-\SubGaugeVec(\CarrierProj\CarrierSeed)_A\bigr)
\bigl(\RelabelSeed^A-\SubGaugeVec(\CarrierProj\CarrierSeed)^A\bigr)
\\
&+\frac{\FreeStiff}{2}\FreeSeed^{AB}\FreeSeed_{AB}
+(\ResidualLift)^{AB}(\CarrierProj\CarrierSeed)_{AB}
\Bigr],
\\
&\qquad\qquad \CarrierGap>0,\qquad \FreeStiff>0.
\end{aligned}
\label{eq:car}
\end{equation}
The primitive readout from that carrier block is
\begin{equation}
\PrimStrain\equiv \CarrierProj\CarrierSeed,
\qquad
\PrimNeutral\equiv \RelabelSeed,
\qquad
\PrimFree\equiv \FreeSeed.
\label{eq:readout}
\end{equation}
\end{definition}

\begin{definition}[Exact-preservation normal-form target]
\label{def:deep}
The precursor-selection deeper-origin note fixed raw deeper readout spaces
\[
\DeepTensorClass,\qquad
\DeepRelabelClass,\qquad
\DeepFreeClass
\]
and bounded surjective linear maps
\begin{equation}
\CarRead:\DeepTensorClass\to\CarrierClass,
\qquad
\CurRead:\DeepRelabelClass\to T^*\Sub,
\qquad
\HomRead:\DeepFreeClass\to\CompFree.
\label{eq:maps}
\end{equation}
For every
\[
\DeepTensor\in\DeepTensorClass,\qquad
\DeepRelabel\in\DeepRelabelClass,\qquad
\DeepFree\in\DeepFreeClass,
\]
write the orthogonal decompositions
\begin{equation}
\DeepTensor=\DeepTensor^{\parallel}+\DeepTensor^{\perp},
\qquad
\DeepTensor^{\parallel}\in(\ker\CarRead)^\perp,
\qquad
\DeepTensor^{\perp}\in\ker\CarRead,
\label{eq:Xsplit}
\end{equation}
\begin{equation}
\DeepRelabel=\DeepRelabel^{\parallel}+\DeepRelabel^{\perp},
\qquad
\DeepRelabel^{\parallel}\in(\ker\CurRead)^\perp,
\qquad
\DeepRelabel^{\perp}\in\ker\CurRead,
\label{eq:Jsplit}
\end{equation}
\begin{equation}
\DeepFree=\DeepFree^{\parallel}+\DeepFree^{\perp},
\qquad
\DeepFree^{\parallel}\in(\ker\HomRead)^\perp,
\qquad
\DeepFree^{\perp}\in\ker\HomRead.
\label{eq:Wsplit}
\end{equation}
Its exact-preservation normal-form functional is
\begin{equation}
\begin{aligned}
\DeepFunctional[\DeepTensor,\DeepRelabel,\DeepFree]
\equiv\;
&\CarrierFunctional[\CarRead\DeepTensor,\CurRead\DeepRelabel,\HomRead\DeepFree]
\\
&+\frac{\PreCarGap}{2}\norm{\DeepTensor^{\perp}}_{\mathrm{deep,car}}^2
+\frac{\CurrentGap}{2}\norm{\DeepRelabel^{\perp}}_{\mathrm{deep,cur}}^2
+\frac{\HomGap}{2}\norm{\DeepFree^{\perp}}_{\mathrm{deep,hom}}^2,
\\
&\qquad\qquad
\PreCarGap>0,\qquad
\CurrentGap>0,\qquad
\HomGap>0.
\end{aligned}
\label{eq:Fdeep}
\end{equation}
\end{definition}

\begin{definition}[Recorded ambient projected-sector construction]
\label{def:proto}
The exact-preservation normal-form origin note fixed ambient substrate spaces
\[
\ProtoTensorClass,\qquad
\ProtoRelabelClass,\qquad
\ProtoFreeClass
\]
containing closed projected sectors identified respectively with
\[
\DeepTensorClass,\qquad
\DeepRelabelClass,\qquad
\DeepFreeClass.
\]
Let
\begin{equation}
\ProtoCarProj:\ProtoTensorClass\to\DeepTensorClass,
\qquad
\ProtoCurProj:\ProtoRelabelClass\to\DeepRelabelClass,
\qquad
\ProtoHomProj:\ProtoFreeClass\to\DeepFreeClass
\label{eq:proto_proj}
\end{equation}
denote the orthogonal projectors onto those closed subspaces, and set
\begin{equation}
\ProtoCarPerp\equiv I-\ProtoCarProj,
\qquad
\ProtoCurPerp\equiv I-\ProtoCurProj,
\qquad
\ProtoHomPerp\equiv I-\ProtoHomProj.
\label{eq:proto_perp}
\end{equation}
Define the ambient visible readout maps by
\begin{equation}
\ProtoCarRead\equiv \CarRead\circ\ProtoCarProj,
\qquad
\ProtoCurRead\equiv \CurRead\circ\ProtoCurProj,
\qquad
\ProtoHomRead\equiv \HomRead\circ\ProtoHomProj.
\label{eq:proto_maps}
\end{equation}
Then
\begin{equation}
\ProtoCarRead\,\ProtoCarPerp=0,
\qquad
\ProtoCurRead\,\ProtoCurPerp=0,
\qquad
\ProtoHomRead\,\ProtoHomPerp=0,
\label{eq:proto_kill}
\end{equation}
and the ambient visible maps have the same codomains
\[
\CarrierClass,\qquad
T^*\Sub,\qquad
\CompFree.
\]
\end{definition}

\begin{definition}[Recorded ambient exact-preservation functional]
\label{def:Fproto}
The previously recorded ambient exact-preservation functional is
\begin{equation}
\begin{aligned}
\ProtoFunctional[\ProtoTensor,\ProtoRelabel,\ProtoFree]
\equiv\;
&\DeepFunctional[
\ProtoCarProj\ProtoTensor,\ProtoCurProj\ProtoRelabel,
\\
&\hspace{1.7em}\ProtoHomProj\ProtoFree]
\\
&+\frac{\ProtoCarGap}{2}\norm{\ProtoCarPerp\ProtoTensor}_{\mathrm{proto,car}}^2
+\frac{\ProtoCurGap}{2}\norm{\ProtoCurPerp\ProtoRelabel}_{\mathrm{proto,cur}}^2
+\frac{\ProtoHomGap}{2}\norm{\ProtoHomPerp\ProtoFree}_{\mathrm{proto,hom}}^2,
\\
&\qquad\qquad
\ProtoCarGap>0,\qquad
\ProtoCurGap>0,\qquad
\ProtoHomGap>0.
\end{aligned}
\label{eq:Fproto}
\end{equation}
\end{definition}

\begin{definition}[Recorded downstream reservoirs]
\label{def:downstream}
The same previous notes already fixed the primitive reservoir
\begin{equation}
\begin{aligned}
\PrimFunctional[\PrimStrain,\PrimNeutral,\PrimFree]
\equiv
\int_{\Sub} d^4X\,\sqrt{G}\,
\Bigl[
&\frac12 \PrimStrain^{AB}\SubEinLin(\PrimStrain)_{AB}
+\frac12\bigl(\PrimNeutral_A-\SubGaugeVec(\PrimStrain)_A\bigr)
\bigl(\PrimNeutral^A-\SubGaugeVec(\PrimStrain)^A\bigr)
\\
&+\frac{\FreeStiff}{2}\PrimFree^{AB}\PrimFree_{AB}
+(\ResidualLift)^{AB}\PrimStrain_{AB}
\Bigr],
\end{aligned}
\label{eq:primitive}
\end{equation}
and the microscopic-equilibrium completion reservoir
\begin{equation}
\begin{aligned}
\MicFunctional[H,\GaugeAux]
=
\int_{\Sub} d^4X\,\sqrt{G}\,
\Bigl[
&\frac12 H^{AB}\SubEinLin(H)_{AB}
+\frac12\bigl(\GaugeAux_A-\SubGaugeVec(H)_A\bigr)
\bigl(\GaugeAux^A-\SubGaugeVec(H)^A\bigr)
\\
&+\frac{\FreeStiff}{2}(\FreeProj H)^{AB}(\FreeProj H)_{AB}
+(\ResidualLift)^{AB}H_{AB}
\Bigr].
\end{aligned}
\label{eq:mic}
\end{equation}
\end{definition}

\begin{definition}[Recorded microphysical exact-preservation functional]
\label{def:Fultra}
The exact-preservation normal-form microphysical-origin note fixed microphysical spaces
\[
\UltraTensorClass,\qquad
\UltraRelabelClass,\qquad
\UltraFreeClass
\]
equipped with defect operators
\begin{equation}
\TensorDefect:\UltraTensorClass\to\mathscr D_{\mathrm{car}},
\qquad
\CurDefect:\UltraRelabelClass\to\mathscr D_{\mathrm{cur}},
\qquad
\HomDefect:\UltraFreeClass\to\mathscr D_{\mathrm{hom}},
\label{eq:defects}
\end{equation}
whose kernels are the exact-preserving ambient sectors
\begin{equation}
\ProtoTensorClass\equiv \ker\TensorDefect,
\qquad
\ProtoRelabelClass\equiv \ker\CurDefect,
\qquad
\ProtoFreeClass\equiv \ker\HomDefect.
\label{eq:proto_kernels}
\end{equation}
Let
\begin{equation}
\UltraCarProj:\UltraTensorClass\to\ProtoTensorClass,
\qquad
\UltraCurProj:\UltraRelabelClass\to\ProtoRelabelClass,
\qquad
\UltraHomProj:\UltraFreeClass\to\ProtoFreeClass
\label{eq:ultra_proj}
\end{equation}
be the orthogonal projectors onto those kernels, and set
\begin{equation}
\UltraCarPerp\equiv I-\UltraCarProj,
\qquad
\UltraCurPerp\equiv I-\UltraCurProj,
\qquad
\UltraHomPerp\equiv I-\UltraHomProj.
\label{eq:ultra_perp}
\end{equation}
Define the visible microphysical maps by
\begin{equation}
\UltraCarRead\equiv \ProtoCarRead\circ\UltraCarProj
=
\CarRead\circ\ProtoCarProj\circ\UltraCarProj,
\label{eq:ultra_car_map}
\end{equation}
\begin{equation}
\UltraCurRead\equiv \ProtoCurRead\circ\UltraCurProj
=
\CurRead\circ\ProtoCurProj\circ\UltraCurProj,
\label{eq:ultra_cur_map}
\end{equation}
\begin{equation}
\UltraHomRead\equiv \ProtoHomRead\circ\UltraHomProj
=
\HomRead\circ\ProtoHomProj\circ\UltraHomProj.
\label{eq:ultra_hom_map}
\end{equation}
The microphysical exact-preservation functional is
\begin{equation}
\begin{aligned}
\UltraFunctional[\UltraTensor,\UltraRelabel,\UltraFree]
\equiv\;
&\ProtoFunctional[
\UltraCarProj\UltraTensor,\UltraCurProj\UltraRelabel,\UltraHomProj\UltraFree]
\\
&+\frac12\norm{\TensorDefect\UltraTensor}_{\mathrm{def,car}}^2
+\frac12\norm{\CurDefect\UltraRelabel}_{\mathrm{def,cur}}^2
+\frac12\norm{\HomDefect\UltraFree}_{\mathrm{def,hom}}^2.
\end{aligned}
\label{eq:Fultra}
\end{equation}
\end{definition}

\begin{remark}
The previous note already established that \eqref{eq:Fultra} reproduces the same ambient projected-sector construction and therefore the same downstream exact-preservation package. The present note does not reopen that ambient or downstream package. It addresses the still sharper question of why the microphysical spaces, the defect operators, and their kernel/gap structure are themselves forced.
\end{remark}

\section{Still-Sharper Symmetry-Compatible Sector}

\begin{definition}[Sharper symmetry projectors and compatibility kernels]
\label{def:root}
Let
\[
\RootTensorClass,\qquad
\RootRelabelClass,\qquad
\RootFreeClass
\]
be Hilbert spaces carrying exact-preserving symmetry projectors
\begin{equation}
\TensorSymProj:\RootTensorClass\to\RootTensorClass,
\qquad
\CurSymProj:\RootRelabelClass\to\RootRelabelClass,
\qquad
\HomSymProj:\RootFreeClass\to\RootFreeClass,
\label{eq:sym_proj}
\end{equation}
with
\[
\TensorSymProj^2=\TensorSymProj=\TensorSymProj^*,
\qquad
\CurSymProj^2=\CurSymProj=\CurSymProj^*,
\qquad
\HomSymProj^2=\HomSymProj=\HomSymProj^*,
\]
and closed compatibility operators
\begin{equation}
\TensorCompat:\operatorname{ran}\TensorSymProj\to\mathscr K_{\mathrm{car}},
\qquad
\CurCompat:\operatorname{ran}\CurSymProj\to\mathscr K_{\mathrm{cur}},
\qquad
\HomCompat:\operatorname{ran}\HomSymProj\to\mathscr K_{\mathrm{hom}}.
\label{eq:compat}
\end{equation}
Define the admissible microphysical sectors by
\begin{equation}
\UltraTensorClass
\equiv
\operatorname{ran}\TensorSymProj\cap\ker\TensorCompat,
\label{eq:ultra_tensor_sector}
\end{equation}
\begin{equation}
\UltraRelabelClass
\equiv
\operatorname{ran}\CurSymProj\cap\ker\CurCompat,
\label{eq:ultra_relabel_sector}
\end{equation}
\begin{equation}
\UltraFreeClass
\equiv
\operatorname{ran}\HomSymProj\cap\ker\HomCompat.
\label{eq:ultra_free_sector}
\end{equation}
Because each space in \eqref{eq:ultra_tensor_sector}--\eqref{eq:ultra_free_sector} is an intersection of closed Hilbert subspaces, there exist orthogonal projectors
\begin{equation}
\RootCarProj:\RootTensorClass\to\UltraTensorClass,
\qquad
\RootCurProj:\RootRelabelClass\to\UltraRelabelClass,
\qquad
\RootHomProj:\RootFreeClass\to\UltraFreeClass
\label{eq:root_proj}
\end{equation}
and orthogonal complements
\begin{equation}
\RootCarPerp\equiv I-\RootCarProj,
\qquad
\RootCurPerp\equiv I-\RootCurProj,
\qquad
\RootHomPerp\equiv I-\RootHomProj.
\label{eq:root_perp}
\end{equation}
\end{definition}

\begin{proposition}[Why the microphysical spaces are forced]
\label{prop:ultra_forced}
If a still sharper substrate layer is to preserve the same three visible codomains
\[
\CarrierClass,\qquad
T^*\Sub,\qquad
\CompFree
\]
and is not allowed to introduce unsuppressed symmetry-breaking or compatibility-breaking exact-preservation data, then the admissible quadratic fluctuations are forced to lie in the closed sectors \eqref{eq:ultra_tensor_sector}--\eqref{eq:ultra_free_sector}.
\end{proposition}

\begin{proof}
Any root fluctuation outside \(\operatorname{ran}\TensorSymProj\), \(\operatorname{ran}\CurSymProj\), or \(\operatorname{ran}\HomSymProj\) carries a symmetry label that is not visible anywhere in the fixed carrier-origin functional \eqref{eq:car}, the fixed ambient projected-sector functional \eqref{eq:Fproto}, or the fixed microphysical functional \eqref{eq:Fultra}. Leaving such directions unsuppressed would therefore enlarge the exact-preservation sector by new symmetry-breaking data. Likewise, any fluctuation inside the symmetry-fixed sector but outside \(\ker\TensorCompat\), \(\ker\CurCompat\), or \(\ker\HomCompat\) violates the compatibility relations required for the same carrier, relabeling-current, and homogeneous slots to survive unchanged. Those directions therefore cannot belong to the admissible unsuppressed exact-preservation sector either.

Hence the only admissible quadratic fluctuations are the simultaneous symmetry-fixed and compatibility-preserving directions \eqref{eq:ultra_tensor_sector}--\eqref{eq:ultra_free_sector}. Since they are closed Hilbert subspaces, the orthogonal projectors \eqref{eq:root_proj} are induced structure rather than same-layer postulates.
\end{proof}

\section{Reduced Equilibrium Hessians and Induced Defect Operators}

\begin{definition}[Reduced equilibrium Hessians]
\label{def:hessians}
Assume the still sharper substrate layer carries symmetry-invariant, compatibility-preserving \(C^2\) response functionals whose exact-preserving equilibrium branch is fixed at the origin of the sectors \(\UltraTensorClass,\UltraRelabelClass,\UltraFreeClass\). Let
\begin{equation}
\TensorHess:\UltraTensorClass\to\UltraTensorClass,
\qquad
\CurHess:\UltraRelabelClass\to\UltraRelabelClass,
\qquad
\HomHess:\UltraFreeClass\to\UltraFreeClass
\label{eq:hess}
\end{equation}
be the nonnegative self-adjoint operators representing the quadratic second variations of that sharper response on the three microphysical sectors. Define the exact-preserving ambient sectors by
\begin{equation}
\ProtoTensorClass\equiv \ker\TensorHess,
\qquad
\ProtoRelabelClass\equiv \ker\CurHess,
\qquad
\ProtoFreeClass\equiv \ker\HomHess.
\label{eq:hess_kernels}
\end{equation}
Let
\begin{equation}
\UltraCarProj:\UltraTensorClass\to\ProtoTensorClass,
\qquad
\UltraCurProj:\UltraRelabelClass\to\ProtoRelabelClass,
\qquad
\UltraHomProj:\UltraFreeClass\to\ProtoFreeClass
\label{eq:hess_proj}
\end{equation}
be the orthogonal projectors onto those kernels, and set
\begin{equation}
\UltraCarPerp\equiv I-\UltraCarProj,
\qquad
\UltraCurPerp\equiv I-\UltraCurProj,
\qquad
\UltraHomPerp\equiv I-\UltraHomProj.
\label{eq:hess_perp}
\end{equation}
\end{definition}

\begin{definition}[Induced defect operators and gap hypothesis]
\label{def:induced_defects}
Identify the defect Hilbert spaces of \eqref{eq:defects} with the Hilbert closures of the ranges of the square roots of the reduced Hessians. Define
\begin{equation}
\TensorDefect\equiv \TensorHess^{1/2},
\qquad
\CurDefect\equiv \CurHess^{1/2},
\qquad
\HomDefect\equiv \HomHess^{1/2}.
\label{eq:defect_sqrt}
\end{equation}
Assume the positive parts of the reduced Hessians satisfy the spectral-gap bounds
\begin{equation}
\langle \TensorHess\UltraTensor,\UltraTensor\rangle_{\mathrm{micro,car}}
\ge
\UltraCarGap\,\norm{\UltraCarPerp\UltraTensor}_{\mathrm{micro,car}}^2,
\qquad
\UltraCarGap>0,
\label{eq:hgap_car}
\end{equation}
\begin{equation}
\langle \CurHess\UltraRelabel,\UltraRelabel\rangle_{\mathrm{micro,cur}}
\ge
\UltraCurGap\,\norm{\UltraCurPerp\UltraRelabel}_{\mathrm{micro,cur}}^2,
\qquad
\UltraCurGap>0,
\label{eq:hgap_cur}
\end{equation}
\begin{equation}
\langle \HomHess\UltraFree,\UltraFree\rangle_{\mathrm{micro,hom}}
\ge
\UltraHomGap\,\norm{\UltraHomPerp\UltraFree}_{\mathrm{micro,hom}}^2,
\qquad
\UltraHomGap>0.
\label{eq:hgap_hom}
\end{equation}
\end{definition}

\begin{proposition}[Why the defect operators are forced]
\label{prop:defects_forced}
At current quadratic exact-preservation scope, the defect operators of the previous note are forced to be the square roots \eqref{eq:defect_sqrt} of the reduced equilibrium Hessians \eqref{eq:hess}.
\end{proposition}

\begin{proof}
Once the admissible microphysical sectors \(\UltraTensorClass,\UltraRelabelClass,\UltraFreeClass\) are fixed by Proposition \ref{prop:ultra_forced}, the sharper substrate response restricts to three nonnegative closed quadratic forms on those Hilbert spaces. By the standard representation theorem for closed nonnegative quadratic forms, each of those forms is represented by a unique nonnegative self-adjoint operator, namely \(\TensorHess,\CurHess,\HomHess\). The spectral theorem then provides unique nonnegative square roots \(\TensorHess^{1/2},\CurHess^{1/2},\HomHess^{1/2}\). Since the exact-preserving ambient sectors are precisely the zero-response directions of the reduced quadratic response, the corresponding defect maps are forced to be those square roots. This is exactly \eqref{eq:defect_sqrt}.
\end{proof}

\begin{proposition}[Closed kernels, orthogonal projectors, and defect-gap estimates]
\label{prop:kernel_gap}
The induced defect operators \eqref{eq:defect_sqrt} recover exactly the kernel/projector/gap package of the previous note:
\begin{enumerate}
    \item \(\ker\TensorDefect=\ProtoTensorClass\), \(\ker\CurDefect=\ProtoRelabelClass\), and \(\ker\HomDefect=\ProtoFreeClass\).
    \item Those kernels are closed Hilbert subspaces, so the orthogonal projectors \(\UltraCarProj,\UltraCurProj,\UltraHomProj\) are forced.
    \item The defect estimates
    \begin{equation}
    \norm{\TensorDefect\UltraTensor}_{\mathrm{def,car}}^2
    \ge
    \UltraCarGap\,\norm{\UltraCarPerp\UltraTensor}_{\mathrm{micro,car}}^2,
    \label{eq:defect_car_recovered}
    \end{equation}
    \begin{equation}
    \norm{\CurDefect\UltraRelabel}_{\mathrm{def,cur}}^2
    \ge
    \UltraCurGap\,\norm{\UltraCurPerp\UltraRelabel}_{\mathrm{micro,cur}}^2,
    \label{eq:defect_cur_recovered}
    \end{equation}
    \begin{equation}
    \norm{\HomDefect\UltraFree}_{\mathrm{def,hom}}^2
    \ge
    \UltraHomGap\,\norm{\UltraHomPerp\UltraFree}_{\mathrm{micro,hom}}^2
    \label{eq:defect_hom_recovered}
    \end{equation}
    hold on the three microphysical sectors.
\end{enumerate}
\end{proposition}

\begin{proof}
Because each Hessian is nonnegative self-adjoint, the square root shares the same kernel. Hence
\[
\ker\TensorDefect=\ker\TensorHess=\ProtoTensorClass,
\qquad
\ker\CurDefect=\ker\CurHess=\ProtoRelabelClass,
\qquad
\ker\HomDefect=\ker\HomHess=\ProtoFreeClass.
\]
The kernels of self-adjoint operators are closed Hilbert subspaces, so the orthogonal projectors \(\UltraCarProj,\UltraCurProj,\UltraHomProj\) exist and are unique.

For the gap estimates, use \(\TensorDefect=\TensorHess^{1/2}\) and the orthogonal decomposition \(\UltraTensor=\UltraCarProj\UltraTensor+\UltraCarPerp\UltraTensor\). Since \(\TensorHess\) vanishes on \(\ker\TensorHess=\operatorname{ran}\UltraCarProj\),
\[
\norm{\TensorDefect\UltraTensor}_{\mathrm{def,car}}^2
=
\langle \TensorHess\UltraTensor,\UltraTensor\rangle_{\mathrm{micro,car}}
=
\langle \TensorHess\UltraCarPerp\UltraTensor,\UltraCarPerp\UltraTensor\rangle_{\mathrm{micro,car}}.
\]
Applying \eqref{eq:hgap_car} gives \eqref{eq:defect_car_recovered}. The proofs of \eqref{eq:defect_cur_recovered} and \eqref{eq:defect_hom_recovered} are identical, using \eqref{eq:hgap_cur} and \eqref{eq:hgap_hom}.
\end{proof}

\section{Exact Recovery of the Same Microphysical and Downstream Package}

\begin{definition}[Reduced symmetry-compatible equilibrium functional]
\label{def:Froot}
The reduced exact-preservation functional induced by the sharper symmetry-compatible equilibrium layer is
\begin{equation}
\begin{aligned}
\RootFunctional[\UltraTensor,\UltraRelabel,\UltraFree]
\equiv\;
&\ProtoFunctional[
\UltraCarProj\UltraTensor,\UltraCurProj\UltraRelabel,\UltraHomProj\UltraFree]
\\
&+\frac12\langle \TensorHess\UltraTensor,\UltraTensor\rangle_{\mathrm{micro,car}}
+\frac12\langle \CurHess\UltraRelabel,\UltraRelabel\rangle_{\mathrm{micro,cur}}
+\frac12\langle \HomHess\UltraFree,\UltraFree\rangle_{\mathrm{micro,hom}}.
\end{aligned}
\label{eq:Froot}
\end{equation}
\end{definition}

\begin{theorem}[Exact recovery of the same microphysical functional]
\label{thm:recover_ultra}
The reduced symmetry-compatible equilibrium functional \eqref{eq:Froot} is exactly the same microphysical functional \eqref{eq:Fultra}:
\begin{equation}
\RootFunctional[\UltraTensor,\UltraRelabel,\UltraFree]
=
\UltraFunctional[\UltraTensor,\UltraRelabel,\UltraFree].
\label{eq:root_equals_ultra}
\end{equation}
\end{theorem}

\begin{proof}
By Proposition \ref{prop:defects_forced},
\[
\TensorDefect=\TensorHess^{1/2},\qquad
\CurDefect=\CurHess^{1/2},\qquad
\HomDefect=\HomHess^{1/2}.
\]
Therefore
\[
\norm{\TensorDefect\UltraTensor}_{\mathrm{def,car}}^2
=
\langle \TensorHess\UltraTensor,\UltraTensor\rangle_{\mathrm{micro,car}},
\]
and similarly for the relabeling and homogeneous sectors. Substituting these identities into \eqref{eq:Fultra} yields exactly \eqref{eq:Froot}. This proves \eqref{eq:root_equals_ultra}.
\end{proof}

\begin{theorem}[Same ambient projected-sector and downstream package]
\label{thm:same_package}
The sharper symmetry-compatible equilibrium layer reproduces exactly the same downstream package already recorded in the exact-preservation normal-form microphysical-origin note:
\begin{enumerate}
    \item the same ambient projected-sector functional \(\ProtoFunctional[\ProtoTensor,\ProtoRelabel,\ProtoFree]\);
    \item the same exact-preservation normal form \(\DeepFunctional[\DeepTensor,\DeepRelabel,\DeepFree]\);
    \item the same carrier-origin functional \(\CarrierFunctional[\CarrierSeed,\RelabelSeed,\FreeSeed]\);
    \item the same primitive reservoir \(\PrimFunctional[\PrimStrain,\PrimNeutral,\PrimFree]\);
    \item the same microscopic-equilibrium completion reservoir \(\MicFunctional[H,\GaugeAux]\);
    \item the same zero-free-mode verdict \(\FreeProj\CompLift=0\);
    \item the same substrate solve \(\SubReducedEin(\CompLift)=-\ResidualLift\);
    \item the same observer solve \(\ReducedEin(\CompShift)=-\ResidualCore\);
    \item the same completion solve \(\Sigma^{\mathrm{cmp}}\), the same one operative metric, and the same recorded Phase D / E and Phase F gains.
\end{enumerate}
\end{theorem}

\begin{proof}
By Theorem \ref{thm:recover_ultra}, the sharper symmetry-compatible equilibrium layer produces exactly the same microphysical functional \eqref{eq:Fultra} as the previous note. The constrained elimination from the microphysical layer to the ambient projected-sector construction, and then onward to the exact-preservation normal form, the carrier-origin block, the primitive reservoir, and the microscopic-equilibrium completion reservoir, depends only on that explicit functional formula and on the recovered defect estimates \eqref{eq:defect_car_recovered}--\eqref{eq:defect_hom_recovered}. Since both are identical to those of the previous note, the stationarity equations and eliminations are unchanged. Every downstream conclusion listed above therefore follows verbatim.
\end{proof}

\begin{corollary}[Recorded gain package is preserved]
\label{cor:gains}
Because the same exact-preservation package is recovered, the recorded Phase D / E infrared-spectrum and universal-coupling verdicts, the recorded Phase F controlled GR limit, and the recorded observer-equation removal of the former genuine quartic residual sector remain preserved without modification.
\end{corollary}

\begin{proof}
This is the direct content of Theorem \ref{thm:same_package}. The sharper note changes only the origin of the microphysical defect construction, not the recovered observer-level completion package.
\end{proof}

\section{Claim-Boundary Verdict}

\begin{theorem}[Microphysical-Hessian origin verdict]
\label{thm:verdict}
At the disciplined scope of the present note, the positive result is exactly the following.
\begin{enumerate}
    \item The microphysical spaces \(\UltraTensorClass,\UltraRelabelClass,\UltraFreeClass\) are no longer left as primitive at current scope.
    \item They are forced to be the closed symmetry-fixed compatibility sectors \eqref{eq:ultra_tensor_sector}--\eqref{eq:ultra_free_sector} of a still sharper substrate layer.
    \item The defect operators \(\TensorDefect,\CurDefect,\HomDefect\) are no longer primitive maps either; they are forced to be the square roots of the reduced equilibrium Hessians \(\TensorHess,\CurHess,\HomHess\).
    \item The exact-preserving ambient sectors \(\ProtoTensorClass,\ProtoRelabelClass,\ProtoFreeClass\) are the zero-response sectors of those reduced Hessians, equivalently the kernels of the induced defect operators.
    \item Because those kernels are closed Hilbert subspaces, the orthogonal kernel projectors \(\UltraCarProj,\UltraCurProj,\UltraHomProj\) are forced.
    \item The positive parts of the reduced Hessians recover exactly the same coercive defect estimates with gaps \(\UltraCarGap,\UltraCurGap,\UltraHomGap\).
    \item The induced reduced symmetry-compatible equilibrium functional is exactly the same microphysical exact-preservation functional \(\UltraFunctional[\UltraTensor,\UltraRelabel,\UltraFree]\) already recorded one layer above.
    \item The same ambient projected-sector construction, the same exact-preservation normal form, the same carrier-origin functional, the same primitive reservoir, the same microscopic-equilibrium reservoir, the same \(\Pi_{\mathrm{free}}H=0\) outcome, the same completion solve \(\Sigma^{\mathrm{cmp}}\), the same observer solve, and the same one operative metric are therefore recovered unchanged.
    \item Stronger promotion language is still not justified here, because the present note explains only the origin of the microphysical defect construction at sharper symmetry-compatible equilibrium scope; it does not yet derive why that still sharper symmetry/compatibility/Hessian package itself is unique, final, or inevitable beyond current scope.
\end{enumerate}
\end{theorem}

\begin{proof}
Items (1) and (2) are Proposition \ref{prop:ultra_forced}. Item (3) is Proposition \ref{prop:defects_forced}. Items (4), (5), and (6) are Proposition \ref{prop:kernel_gap}. Item (7) is Theorem \ref{thm:recover_ultra}. Item (8) is Theorem \ref{thm:same_package} together with Corollary \ref{cor:gains}. Item (9) is the claim-boundary discipline stated in the introduction.
\end{proof}

\begin{corollary}[What remains next]
The primary active burden moves one layer deeper again. It is no longer to justify why the defect operators, their closed defect-null kernels, their orthogonal kernel projectors, and their coercive positive-gap estimates are admissible at current scope; that step is now recorded. The next honest burden is either to derive or justify why the still sharper symmetry projectors, compatibility kernels, and gapped reduced equilibrium Hessians themselves arise from yet more primitive substrate structure, or else to produce another comparably concrete observer-equation gain on the same one-metric line. The anchored positive-pair continuation remains side work unless it directly supplies one of those.
\end{corollary}

\section{Conclusion}

The positive result of this note is that the defect-null microphysical construction is no longer left unexplained at present scope. It is recovered from a sharper symmetry-compatible equilibrium-response layer.

At that sharper level, one begins with still sharper substrate spaces
\[
\RootTensorClass,\qquad
\RootRelabelClass,\qquad
\RootFreeClass
\]
equipped with exact-preserving symmetry projectors
\[
\TensorSymProj,\qquad
\CurSymProj,\qquad
\HomSymProj
\]
and closed compatibility operators
\[
\TensorCompat,\qquad
\CurCompat,\qquad
\HomCompat.
\]
The admissible microphysical sectors are then the closed symmetry-fixed compatibility sectors
\[
\UltraTensorClass=\operatorname{ran}\TensorSymProj\cap\ker\TensorCompat,
\qquad
\UltraRelabelClass=\operatorname{ran}\CurSymProj\cap\ker\CurCompat,
\qquad
\UltraFreeClass=\operatorname{ran}\HomSymProj\cap\ker\HomCompat.
\]
On those sectors, the sharper equilibrium response has nonnegative reduced Hessians
\[
\TensorHess,\qquad
\CurHess,\qquad
\HomHess.
\]
Their square roots are exactly the defect operators
\[
\TensorDefect=\TensorHess^{1/2},
\qquad
\CurDefect=\CurHess^{1/2},
\qquad
\HomDefect=\HomHess^{1/2}.
\]
The exact-preserving ambient sectors are therefore the zero-response directions
\[
\ProtoTensorClass=\ker\TensorHess=\ker\TensorDefect,
\qquad
\ProtoRelabelClass=\ker\CurHess=\ker\CurDefect,
\qquad
\ProtoFreeClass=\ker\HomHess=\ker\HomDefect.
\]
Because those kernels are closed Hilbert subspaces, the orthogonal kernel projectors are induced automatically, and the positive parts of the Hessians recover the same coercive defect-gap estimates with gaps \(\UltraCarGap,\UltraCurGap,\UltraHomGap\).

The reduced symmetry-compatible equilibrium functional is therefore
\[
\RootFunctional[\UltraTensor,\UltraRelabel,\UltraFree],
\]
and it is exactly the same microphysical functional already recorded in the previous note:
\[
\RootFunctional[\UltraTensor,\UltraRelabel,\UltraFree]
=
\UltraFunctional[\UltraTensor,\UltraRelabel,\UltraFree].
\]
Hence the same ambient projected-sector construction, the same exact-preservation normal form, the same carrier-origin functional, the same primitive reservoir, the same microscopic-equilibrium reservoir, the same \(\Pi_{\mathrm{free}}H=0\) outcome, the same substrate and observer solves, the same completion \(\Sigma^{\mathrm{cmp}}\), and the same one operative metric survive unchanged. The recorded Phase D / E and Phase F gains remain intact.

The scope remains honest. This note does not claim a final ultraviolet completion and does not widen the current theorem boundary. Its exact positive result is narrower and precisely the one now needed: the previous defect-null microphysical construction is no longer primitive at current scope, but is forced by a sharper symmetry-compatible equilibrium-response structure while preserving the same downstream package. The next honest burden is one layer deeper again: whether that still sharper symmetry/compatibility/Hessian package can itself be derived or selected from yet more primitive substrate structure, or else superseded by another equally concrete observer-equation gain on the same one-metric line.

\end{document}
